% Supplementary Note 7: what the explanations referred to. Numbers through macros only.
\suppnote{What the explanations referred to}

This note gives in full the analyses of participants' written explanations that the main text summarizes in its Results section of the same name.

\paragraph*{How the explanations were read.}
In all five studies, a required text box for each rating asked participants to explain their reasoning, and every box came after the ratings it explained. Eight fixed word lists (for example, words for visible detail, time, skill, and prompting or machines) flagged what each explanation mentioned. We wrote the lists by reading Study~1's explanations with their label and rating hidden, and fixed them before relating them to either, apart from four parts settled after results by label had been read (Section~\ref{sec:suppA-text}). A script then applied the lists without seeing either, and no person or language model classified any single explanation. Every explanation was written after its rating, so these results describe what the words referred to, not why anyone rated as they did. Everything in this note is exploratory, meaning that a registration named it as exploratory or did not name it at all, or descriptive, meaning a share or a count with no test. The analyses of second-work explanations in Studies 2, 4 and 5, of talk of origin and of how far explanations moved, were specified on 9 October 2026, after every registered result was known, and they change no registered verdict. A difference between two shares is given in percentage points, the plain difference between two percentages.

\paragraph*{Originals.}
On originals, what `effort' referred to depended on the label. In Study~1, \SOneN{} participants rated ten AI-generated images on creativity and on effort, told either that Alex made them with design and sketching tools (the Human label) or that he made them by prompting an AI system (the AI label). Study~1's registration named five word-meaning tests as one exploratory family and set the Benjamini--Hochberg correction for them. That correction limits the expected share of false positives among results called significant;\cite{benjamini1995controlling} a result is called significant when its $p$-value, the chance of a difference at least this large if there were none, is below .05. The corrected results are descriptive readings, not verdicts. Each test compared how often two kinds of mention occurred, allowing for differences between participants and between images. Under the Human label, \SOneTextEffortDetailPct{} of effort explanations named visible detail, \SOneTextEffortTimePct{} named time and \SOneTextEffortSkillPct{} named skill. Detail was named more often than time and more often than skill (Benjamini--Hochberg $p$~\SOneTextDetailVsTimePBH{} against time and $p$~\SOneTextDetailVsSkillPBH{} against skill). Under the AI label, \SOneTextPromptEffortAIPct{} of effort explanations mentioned prompting or machines, against \SOneTextPromptEffortHumanPct{} under the Human label (descriptive). The difference held among explanations given with the same effort rating of 3 (\SOneTextPromptEffortRatingThreeAIPct{} against \SOneTextPromptEffortRatingThreeHumanPct; descriptive), so it does not arise only because AI-label raters gave lower ratings. In a third test of the family, AI-label participants mentioned prompting or machines more often when they explained effort than when they explained creativity (Benjamini--Hochberg $p$~\SOneTextPromptEffortVsCreatPBH). The family's last two tests concerned Human-label creativity explanations. These named the idea behind the work (\SOneTextCreatConceptPct) more often than its execution (\SOneTextCreatExecutionPct; Benjamini--Hochberg $p$~\SOneTextConceptVsExecutionPBH), but not detectably more often than whether it looked new or familiar (\SOneTextCreatNoveltyPct; $p$~\SOneTextConceptVsNoveltyPBH). How the word list for new or familiar handles negations was settled after results by label had been read. Because these two tests used Human-label explanations only, they say nothing about what AI-label creativity ratings rested on.

Study~3, in which \SThreeN{} participants rated nine of the originals under the Human or the AI label, each with a stated time of under 1, 10 or 100 hours, repeated the pattern on effort. Each difference below is AI minus Human, in percentage points, with a 95\% confidence interval (CI), the range of differences the data do not rule out at the 95\% level. AI-label effort explanations named detail less often and prompting more often (\SThreeTextWordEffDetailDiff{} for detail, 95\% CI \SThreeTextWordEffDetailLo{} to \SThreeTextWordEffDetailHi; \SThreeTextWordEffPromptDiff{} for prompting, \SThreeTextWordEffPromptLo{} to \SThreeTextWordEffPromptHi). AI-label creativity explanations also mentioned prompting more often (\SThreeTextWordCreatPromptDiff{} points, \SThreeTextWordCreatPromptLo{} to \SThreeTextWordCreatPromptHi). These intervals come from resampling participants, and none of Study~3's word-list comparisons was corrected for multiple comparisons, so an interval that only just excludes zero is tentative. Rating the same pictures, the two label groups wrote about different acts when they explained effort. This does not show that detail drove effort ratings or that `effort' meant labor.

\paragraph*{Second works: talk of origin.}
On second works, the label left a faint trace in what people wrote. In Studies 2, 4 and 5, participants explained their ratings of second works that Alex was said to have made from each original with his design and sketching tools: None, the original shown again unchanged; Slight, a light edit; and Dramatic, a heavy one. Both label groups were told that Alex made every second work this way; only the description of the original differed. Pooled over the three studies, with each study weighted equally, AI-label explanations mentioned prompting or a machine more often than Human-label ones at every level (\XSOriginTalkAINonePct, \XSOriginTalkAISlightPct{} and \XSOriginTalkAIDramaticPct{} at None, Slight and Dramatic, against \XSOriginTalkHumanNonePct, \XSOriginTalkHumanSlightPct{} and \XSOriginTalkHumanDramaticPct{} under the Human label). These shares count explanations of all three ratings from people who explained every level of the same original, so each level holds the same people and images. Such mentions rose from None to Dramatic under both labels, and the two rises did not differ detectably (difference \XSOriginTalkChangeDiffPts{} percentage points, 95\% CI \XSOriginTalkChangeDiffLoPts{} to \XSOriginTalkChangeDiffHiPts, from resampling participants). In Study~2, however, AI-label creativity explanations went from \XSOriginTalkSTwoCreatAINonePct{} at None to \XSOriginTalkSTwoCreatAIDramaticPct{} at Dramatic, against \XSOriginTalkSTwoCreatHumanNonePct{} and \XSOriginTalkSTwoCreatHumanDramaticPct{} under the Human label, so their rise was \XSOriginTalkSTwoCreatChangeDiffPts{} percentage points larger (95\% CI \XSOriginTalkSTwoCreatChangeDiffLoPts{} to \XSOriginTalkSTwoCreatChangeDiffHiPts; uncorrected).

\paragraph*{Second works: how far explanations moved.}
Under taint, the account on which the AI origin passes to anything built from it, AI-label raters should have kept writing about the original, so their explanations should have changed less from None to Dramatic. That change was measured in two ways: with sentence embeddings, which turn each explanation into a list of numbers whose distances reflect differences in meaning, and with word overlap, the share of content words two explanations have in common.\cite{reimers2019sentencebert,sentencetransformers2021allmpnet} For each person and image, the measure began with how far that person's Dramatic explanation lay from their None explanation. From this it subtracted how far their None explanation lay from other participants' Dramatic explanations of the same image, which removes what any two people writing about the same picture share. The label groups were then compared with explanation length held constant; a negative difference, AI minus Human, means that AI-label explanations moved less. Pooled over studies, no difference in change was shown between the label groups. On creativity the difference was \XSCreatShiftEmbEst{} with sentence embeddings (95\% CI \XSCreatShiftEmbLo{} to \XSCreatShiftEmbHi) and \XSCreatShiftWordEst{} with word overlap (\XSCreatShiftWordLo{} to \XSCreatShiftWordHi). The pooled estimates all lay on the side taint predicts (\XSShiftNPooledTaintLean{} of \XSShiftNPooled, two distances on each of three measures), but none was clear of zero. Of \XSShiftNRows{} uncorrected comparisons, pooled and study by study, \XSShiftNRowsTaint{} fell clear of zero on that side (Study~5's transformativeness explanations, word overlap), about what chance alone would produce. No difference shown is not the same as no difference. No study asked participants to recall the label (a manipulation check), so none of this shows that the label was in mind when people rated.
